\documentclass[9pt,a4paper]{extarticle}

% ==========================================
% Basic packages
% ==========================================
\usepackage[T1]{fontenc}
\usepackage[utf8]{inputenc}
\usepackage[brazil]{babel}
\usepackage{geometry}
\geometry{margin=2cm}
\usepackage{parskip}
\usepackage{microtype}

% Math
\usepackage{amsmath, amssymb, amsthm, bm}

% ==========================================
% Document
% ==========================================
\begin{document}

\section*{Lecture Notes: Fletcher--Reeves Method (Conjugate Gradient)}

\section{Motivation}

The Conjugate Gradient method is an efficient alternative to steepest descent for minimizing quadratic functions of the form
\[
f(x)=\tfrac12 x^\top A x + b^\top x + c,
\]
where $A$ is symmetric positive definite. Unlike steepest descent, which may zig-zag, the Conjugate Gradient method constructs mutually $A$-conjugate directions, ensuring convergence in at most $n$ iterations for quadratic functions.

\section{Line Minimization}

Consider a quadratic objective function of the form
\[
f(x) = \tfrac{1}{2} x^\top A x + b^\top x + c,
\]
where $A$ is symmetric (and, no contexto do método, definida positiva).

Given a current point $x_i$ and a search direction $S_i$, we define the next iterate as
\[
x_{i+1} = x_i + \lambda_i S_i,
\]
where the scalar $\lambda_i$ is chosen to minimize $f$ along the line
\[
\phi(\lambda) = f(x_i + \lambda S_i).
\]

\subsection*{Step 1: Gradient of the quadratic function}

For the quadratic function above, the gradient is
\[
\nabla f(x) = A x + b.
\]

Em particular, no ponto $x_i$ temos
\[
\nabla f_i := \nabla f(x_i) = A x_i + b.
\]

\subsection*{Step 2: Expressing the gradient along the line}

Ao longo da reta $x(\lambda) = x_i + \lambda S_i$, o gradiente em função de $\lambda$ é
\[
\nabla f(x_i + \lambda S_i)
= A(x_i + \lambda S_i) + b.
\]

Usando a distributividade,
\[
\nabla f(x_i + \lambda S_i)
= A x_i + \lambda A S_i + b.
\]

Reagrupando os termos conhecidos em $x_i$:
\[
\nabla f(x_i + \lambda S_i)
= (A x_i + b) + \lambda A S_i
= \nabla f_i + \lambda A S_i.
\]

\subsection*{Step 3: Derivative of the line function \texorpdfstring{$\phi(\lambda)$}{phi(lambda)}}

A função de uma variável real que estamos minimizando é
\[
\phi(\lambda) = f(x_i + \lambda S_i).
\]

Pela regra da cadeia,
\[
\frac{d\phi}{d\lambda}
= \nabla f(x_i + \lambda S_i)^\top \frac{d}{d\lambda}(x_i + \lambda S_i).
\]

Como
\[
\frac{d}{d\lambda}(x_i + \lambda S_i) = S_i,
\]
temos
\[
\frac{d\phi}{d\lambda}
= \nabla f(x_i + \lambda S_i)^\top S_i.
\]

Substituindo a expressão da Subsecção 2:
\[
\frac{d\phi}{d\lambda}
= (\nabla f_i + \lambda A S_i)^\top S_i.
\]

\subsection*{Step 4: Simplifying the derivative}

Usando $(u+v)^\top S_i = u^\top S_i + v^\top S_i$:
\[
\frac{d\phi}{d\lambda}
= \nabla f_i^\top S_i + (\lambda A S_i)^\top S_i.
\]

Note que $(\lambda A S_i)^\top S_i = \lambda (A S_i)^\top S_i$.  
Como $A$ é simétrica, $(A S_i)^\top S_i = S_i^\top A S_i$. Logo:
\[
\frac{d\phi}{d\lambda}
= \nabla f_i^\top S_i + \lambda S_i^\top A S_i.
\]

Podemos também escrever
\[
\frac{d\phi}{d\lambda}
= S_i^\top \nabla f_i + \lambda S_i^\top A S_i.
\]

\subsection*{Step 5: Optimal step size}

O valor ótimo $\lambda_i^*$ é obtido impondo a condição de primeira ordem
\[
\frac{d\phi}{d\lambda} = 0.
\]

Assim,
\[
S_i^\top \nabla f_i + \lambda_i^* S_i^\top A S_i = 0.
\]

Isolando $\lambda_i^*$:
\[
\lambda_i^* S_i^\top A S_i = - S_i^\top \nabla f_i,
\]
\[
\lambda_i^* = -\frac{S_i^\top \nabla f_i}{S_i^\top A S_i}.
\]

Em particular, na primeira iteração ($i=1$) temos
\[
\lambda_1^* = -\frac{S_1^\top \nabla f_1}{S_1^\top A S_1}.
\]


\section{Conjugate Directions}

We now show in detail how to construct the second search direction $S_2$ and how to obtain
the scalar $\beta_2$ that makes $S_1$ and $S_2$ $A$-conjugate.

\subsection*{Step 1: First search direction}

The first search direction is chosen as the negative gradient:
\[
S_1 = - \nabla f_1,
\]
where
\[
\nabla f_1 := \nabla f(x_1).
\]

\subsection*{Step 2: Second search direction as a combination}

We assume that the second direction is a linear combination of the new steepest descent
direction $-\nabla f_2$ and the previous direction $S_1$:
\[
S_2 = -\nabla f_2 + \beta_2 S_1,
\]
where
\[
\nabla f_2 := \nabla f(x_2),
\qquad
x_2 = x_1 + \lambda_1^* S_1,
\]
and $\beta_2$ is an unknown scalar to be determined.

\subsection*{Step 3: Conjugacy condition}

We want $S_1$ and $S_2$ to be $A$-conjugate, i.e.,
\[
S_1^\top A S_2 = 0.
\]

Substitute the expression of $S_2$:
\[
S_1^\top A(-\nabla f_2 + \beta_2 S_1) = 0.
\]

Distribute $A$ and separate the terms:
\[
- S_1^\top A \nabla f_2 + \beta_2 S_1^\top A S_1 = 0.
\]

Now solve for $\beta_2$:
\[
\beta_2 S_1^\top A S_1 = S_1^\top A \nabla f_2,
\]
\[
\beta_2 = \frac{S_1^\top A \nabla f_2}{S_1^\top A S_1}.
\]

Neste ponto, a expressão ainda contém a matriz $A$. O próximo passo é reescrevê-la usando apenas gradientes.

\subsection*{Step 4: Relation between consecutive gradients}

Recall that, for the quadratic function
\[
f(x) = \tfrac{1}{2} x^\top A x + b^\top x + c,
\]
the gradient is
\[
\nabla f(x) = A x + b.
\]

At $x_1$ and $x_2$:
\[
\nabla f_1 = A x_1 + b, 
\qquad
\nabla f_2 = A x_2 + b.
\]

Since
\[
x_2 = x_1 + \lambda_1^* S_1,
\]
we have
\[
\nabla f_2
= A(x_1 + \lambda_1^* S_1) + b
= A x_1 + \lambda_1^* A S_1 + b
= (A x_1 + b) + \lambda_1^* A S_1
= \nabla f_1 + \lambda_1^* A S_1.
\]

Rearranging, we obtain:
\[
A S_1 = \frac{\nabla f_2 - \nabla f_1}{\lambda_1^*}.
\]

\subsection*{Step 5: Denominator $S_1^\top A S_1$ in terms of gradients}

Use the identity above:
\[
S_1^\top A S_1
= S_1^\top \left( \frac{\nabla f_2 - \nabla f_1}{\lambda_1^*} \right)
= \frac{S_1^\top \nabla f_2 - S_1^\top \nabla f_1}{\lambda_1^*}.
\]

Next, we will use the fact that the line search is exact.

\subsection*{Step 6: Orthogonality from exact line search}

Consider the line function
\[
\phi(\lambda) = f(x_1 + \lambda S_1).
\]

We know that the derivative at the optimal step $\lambda_1^*$ is zero:
\[
\left. \frac{d\phi}{d\lambda} \right|_{\lambda = \lambda_1^*}
= 0.
\]

As seen before,
\[
\frac{d\phi}{d\lambda}
= \nabla f(x_1 + \lambda S_1)^\top S_1,
\]
thus, at $\lambda = \lambda_1^*$,
\[
\nabla f_2^\top S_1 = 0
\quad\Longrightarrow\quad
S_1^\top \nabla f_2 = 0.
\]

Substitute this into the denominator:
\[
S_1^\top A S_1
= \frac{0 - S_1^\top \nabla f_1}{\lambda_1^*}
= - \frac{S_1^\top \nabla f_1}{\lambda_1^*}.
\]

Using $S_1 = -\nabla f_1$:
\[
S_1^\top \nabla f_1 = - \nabla f_1^\top \nabla f_1,
\]
so
\[
S_1^\top A S_1
= - \frac{- \nabla f_1^\top \nabla f_1}{\lambda_1^*}
= \frac{\nabla f_1^\top \nabla f_1}{\lambda_1^*}.
\]

\subsection*{Step 7: Numerator $S_1^\top A \nabla f_2$ in terms of gradients}

Using $A S_1 = (\nabla f_2 - \nabla f_1)/\lambda_1^*$:
\[
S_1^\top A \nabla f_2
= (A S_1)^\top \nabla f_2
= \left( \frac{\nabla f_2 - \nabla f_1}{\lambda_1^*} \right)^\top \nabla f_2.
\]

Compute the product:
\[
S_1^\top A \nabla f_2
= \frac{\nabla f_2^\top \nabla f_2 - \nabla f_1^\top \nabla f_2}{\lambda_1^*}.
\]

Agora usamos novamente uma ortogonalidade importante.

Note that
\[
S_1 = -\nabla f_1
\quad\Longrightarrow\quad
S_1^\top \nabla f_2 = - \nabla f_1^\top \nabla f_2.
\]

But from the exact line search condition we already know that
\[
S_1^\top \nabla f_2 = 0
\quad\Longrightarrow\quad
\nabla f_1^\top \nabla f_2 = 0.
\]

Thus,
\[
S_1^\top A \nabla f_2
= \frac{\nabla f_2^\top \nabla f_2 - 0}{\lambda_1^*}
= \frac{\nabla f_2^\top \nabla f_2}{\lambda_1^*}.
\]

\subsection*{Step 8: Final expression for \texorpdfstring{$\beta_2$}{beta2}}

Recall
\[
\beta_2 = \frac{S_1^\top A \nabla f_2}{S_1^\top A S_1}.
\]

Substitute numerator and denominator:
\[
\beta_2
= \frac{\dfrac{\nabla f_2^\top \nabla f_2}{\lambda_1^*}}{\dfrac{\nabla f_1^\top \nabla f_1}{\lambda_1^*}}
= \frac{\nabla f_2^\top \nabla f_2}{\nabla f_1^\top \nabla f_1}.
\]

This is exactly the Fletcher--Reeves formula for $\beta_2$.

\section{General Formula for \texorpdfstring{$\beta_i$}{beta\_i}}

We now generalize the previous derivation to obtain the expression for $\beta_i$ for $i \ge 2$.

\subsection*{Step 1: General search direction}

Assume that, at iteration $i$, the search direction is defined by
\[
S_i = -\nabla f_i + \beta_i S_{i-1},
\]
where $\nabla f_i := \nabla f(x_i)$.

We require $S_{i-1}$ and $S_i$ to be $A$-conjugate:
\[
S_{i-1}^\top A S_i = 0.
\]

Substitute $S_i$:
\[
S_{i-1}^\top A(-\nabla f_i + \beta_i S_{i-1}) = 0.
\]

Expand:
\[
- S_{i-1}^\top A \nabla f_i + \beta_i S_{i-1}^\top A S_{i-1} = 0.
\]

Solve for $\beta_i$:
\[
\beta_i S_{i-1}^\top A S_{i-1} = S_{i-1}^\top A \nabla f_i,
\]
\[
\beta_i = \frac{S_{i-1}^\top A \nabla f_i}{S_{i-1}^\top A S_{i-1}}.
\]

Again, we want to eliminate $A$ using gradient relations.

\subsection*{Step 2: Relation between \texorpdfstring{$\nabla f_i$}{grad fi} and \texorpdfstring{$\nabla f_{i-1}$}{grad f(i-1)}}

As before, for the quadratic function
\[
\nabla f(x) = A x + b.
\]

At iterations $i-1$ and $i$:
\[
\nabla f_{i-1} = A x_{i-1} + b,
\qquad
\nabla f_i = A x_i + b.
\]

We also have
\[
x_i = x_{i-1} + \lambda_{i-1}^* S_{i-1}.
\]

Therefore,
\[
\nabla f_i
= A(x_{i-1} + \lambda_{i-1}^* S_{i-1}) + b
= (A x_{i-1} + b) + \lambda_{i-1}^* A S_{i-1}
= \nabla f_{i-1} + \lambda_{i-1}^* A S_{i-1}.
\]

Rearranging:
\[
A S_{i-1} = \frac{\nabla f_i - \nabla f_{i-1}}{\lambda_{i-1}^*}.
\]

\subsection*{Step 3: Denominator $S_{i-1}^\top A S_{i-1}$}

Compute:
\[
S_{i-1}^\top A S_{i-1}
= S_{i-1}^\top \left( \frac{\nabla f_i - \nabla f_{i-1}}{\lambda_{i-1}^*} \right)
= \frac{S_{i-1}^\top \nabla f_i - S_{i-1}^\top \nabla f_{i-1}}{\lambda_{i-1}^*}.
\]

As in the case $i=1$, the exact line search at step $i-1$ implies
\[
\left. \frac{d}{d\lambda} f(x_{i-1} + \lambda S_{i-1}) \right|_{\lambda = \lambda_{i-1}^*}
= 0,
\]
which gives
\[
\nabla f_i^\top S_{i-1} = 0
\quad\Longrightarrow\quad
S_{i-1}^\top \nabla f_i = 0.
\]

Thus,
\[
S_{i-1}^\top A S_{i-1}
= \frac{0 - S_{i-1}^\top \nabla f_{i-1}}{\lambda_{i-1}^*}
= - \frac{S_{i-1}^\top \nabla f_{i-1}}{\lambda_{i-1}^*}.
\]

Using the structure of $S_{i-1}$, one can show that the gradients are mutually orthogonal and, in particular,
\[
\nabla f_{i-1}^\top \nabla f_i = 0,
\]
and $S_{i-1}$ has a component $-\nabla f_{i-1}$ dominating the product with $\nabla f_{i-1}$. Em termos práticos, obtém-se
\[
S_{i-1}^\top A S_{i-1} = \frac{\nabla f_{i-1}^\top \nabla f_{i-1}}{\lambda_{i-1}^*}.
\]

(Para uma prova completa, segue-se o mesmo raciocínio recursivo do caso $i=1$.)

\subsection*{Step 4: Numerator $S_{i-1}^\top A \nabla f_i$}

Using $A S_{i-1} = (\nabla f_i - \nabla f_{i-1})/\lambda_{i-1}^*$:
\[
S_{i-1}^\top A \nabla f_i
= (A S_{i-1})^\top \nabla f_i
= \left( \frac{\nabla f_i - \nabla f_{i-1}}{\lambda_{i-1}^*} \right)^\top \nabla f_i.
\]

Compute:
\[
S_{i-1}^\top A \nabla f_i
= \frac{\nabla f_i^\top \nabla f_i - \nabla f_{i-1}^\top \nabla f_i}{\lambda_{i-1}^*}.
\]

Again, by orthogonality of consecutive gradients (from exact line search),
\[
\nabla f_{i-1}^\top \nabla f_i = 0,
\]
so
\[
S_{i-1}^\top A \nabla f_i
= \frac{\nabla f_i^\top \nabla f_i}{\lambda_{i-1}^*}.
\]

\subsection*{Step 5: Final expression for \texorpdfstring{$\beta_i$}{beta\_i}}

Now,
\[
\beta_i = \frac{S_{i-1}^\top A \nabla f_i}{S_{i-1}^\top A S_{i-1}}
= \frac{\dfrac{\nabla f_i^\top \nabla f_i}{\lambda_{i-1}^*}}{\dfrac{\nabla f_{i-1}^\top \nabla f_{i-1}}{\lambda_{i-1}^*}}
= \frac{\nabla f_i^\top \nabla f_i}{\nabla f_{i-1}^\top \nabla f_{i-1}}.
\]

Thus, for all $i \ge 2$,
\[
\beta_i =
\frac{\nabla f_i^\top \nabla f_i}{\nabla f_{i-1}^\top \nabla f_{i-1}}.
\]

This is the Fletcher--Reeves update formula used in the Conjugate Gradient method.

\section{Fletcher--Reeves Algorithm (English Version)}

\begin{enumerate}
    \item Choose initial guess $x_1$.
    \item Compute $\nabla f_1$ and set $S_1 = -\nabla f_1$.
    \item For $i = 1, 2, \dots$:
    \begin{enumerate}
        \item Compute the step size
        \[
        \lambda_i^* = -\frac{S_i^\top \nabla f_i}{S_i^\top A S_i}.
        \]
        \item Update the iterate
        \[
        x_{i+1} = x_i + \lambda_i^* S_i.
        \]
        \item Compute the new gradient $\nabla f_{i+1}$.
        \item If $\|\nabla f_{i+1}\|$ is sufficiently small, stop.
        \item Compute the scalar
        \[
        \beta_{i+1} = \frac{\nabla f_{i+1}^\top \nabla f_{i+1}}{\nabla f_i^\top \nabla f_i}.
        \]
        \item Update the direction
        \[
        S_{i+1} = -\nabla f_{i+1} + \beta_{i+1} S_i.
        \]
    \end{enumerate}
\end{enumerate}

\section{Example 6.9}

Minimize:
\[
f(x_1, x_2) = x_1 - x_2 + 2x_1^2 + 2x_1 x_2 + x_2^2,
\]
starting from:
\[
X_1 =
\begin{bmatrix}
0\\
0
\end{bmatrix}.
\]

\subsection*{Iteration 1}

The gradient is:
\[
\nabla f =
\begin{bmatrix}
1 + 4x_1 + 2x_2 \\
-1 + 2x_1 + 2x_2
\end{bmatrix}.
\]

Thus:
\[
\nabla f_1 = \nabla f(0,0) =
\begin{bmatrix}
1\\
-1
\end{bmatrix},
\qquad
S_1 = -\nabla f_1 =
\begin{bmatrix}
-1\\
1
\end{bmatrix}.
\]

Along this direction:
\[
f(X_1 + \lambda_1 S_1)
= f(-\lambda_1, +\lambda_1) = \lambda_1^2 - 2\lambda_1.
\]

Derivative:
\[
\frac{d}{d\lambda_1} = 2\lambda_1 - 2 = 0
\quad\Rightarrow\quad
\lambda_1^* = 1.
\]

Update:
\[
X_2 = X_1 + S_1 =
\begin{bmatrix}
-1\\
1
\end{bmatrix}.
\]

\subsection*{Iteration 2}

\[
\nabla f_2 = \nabla f(-1,1) =
\begin{bmatrix}
-1\\
-1
\end{bmatrix}.
\]

Norms:
\[
\|\nabla f_1\|^2 = 2,
\qquad
\|\nabla f_2\|^2 = 2.
\]

Direction:
\[
S_2 = -\nabla f_2 + \frac{\|\nabla f_2\|^2}{\|\nabla f_1\|^2} S_1
=
\begin{bmatrix}
0\\
2
\end{bmatrix}.
\]

Compute step:
\[
f(X_2 + \lambda_2 S_2)
= f(-1,1 + 2\lambda_2)
= 4\lambda_2^2 - 2\lambda_2 - 1.
\]

Derivative:
\[
8\lambda_2 - 2 = 0
\Rightarrow
\lambda_2^* = \frac14.
\]

Update:
\[
X_3 =
\begin{bmatrix}
-1\\[4pt]
1.5
\end{bmatrix}.
\]

\subsection*{Iteration 3}

\[
\nabla f_3 = \nabla f(-1,1.5)=
\begin{bmatrix}
0\\
0
\end{bmatrix}.
\]

Thus:
\[
\beta_3 = 0, 
\qquad
S_3 = 0.
\]

No further descent direction exists, hence $X_3$ is the minimizer.

\section{Conclusion}

The Fletcher--Reeves method constructs conjugate directions using only gradient information. For quadratic functions, it converges in at most $n$ iterations, as illustrated in the example.

\section{Fraquezas do Método de Fletcher--Reeves}

Embora o método de Fletcher--Reeves (FR) seja um dos formuladores clássicos do Método dos Gradientes Conjugados, seu desempenho prático apresenta limitações importantes. A seguir destacamos suas principais fraquezas, especialmente relevantes em problemas não quadráticos ou mal condicionados.

\subsection*{Sensibilidade ao \textit{line search}}

O método FR depende criticamente de uma busca linear que satisfaça as condições de Wolfe de forma relativamente precisa. Quando o \textit{line search} é inexato ou instável, o valor de $\beta_i$ tende a oscilar, e as direções deixam de ser verdadeiramente $A$-conjugadas. Como consequência, o algoritmo perde suas propriedades teóricas e passa a se comportar de modo semelhante ao método de gradiente descendente.

\subsection*{Perda de conjugação devido a erros numéricos}

A construção das direções conjugadas pressupõe que operações envolvendo o gradiente e a matriz $A$ (implícita) sejam numericamente estáveis. Em problemas de grande dimensão ou com forte assimetria de escalas, erros de arredondamento acumulam-se ao longo das iterações, destruindo a relação
\[
S_i^\top A S_j = 0 \quad (i \neq j).
\]
Quando isso ocorre, o método perde eficiência e necessita de reinicializações frequentes.

\subsection*{Desempenho ruim em problemas mal condicionados}

Quando a matriz Hessiana (ou $A$, no caso quadrático) é mal condicionada, isto é, possui número de condição elevado,
\[
\kappa(A) = \frac{\lambda_{\max}}{\lambda_{\min}} \gg 1,
\]
as curvas de nível tornam-se elipses muito alongadas. Nesse cenário, os gradientes mudam de direção abruptamente, levando a valores de $\beta_i$ pouco informativos e dificultando a manutenção das direções conjugadas. Na prática, o método tende a gerar trajetórias em ``zigue-zague'' e convergência lenta.

\subsection*{Fragilidade em funções não quadráticas}

A teoria que garante convergência em até $n$ passos vale apenas para funções quadráticas. Em funções gerais (como as encontradas em otimização não linear ou aprendizado de máquina), a relação exata entre gradientes consecutivos deixa de valer, comprometendo a validade de $\beta_i^{FR}$. O método pode apresentar estagnação, oscilações ou necessidade de reinicializações constantes.

\subsection*{Necessidade de reinicializações periódicas}

Para mitigar perda de conjugação, é comum reiniciar o método definindo
\[
S_i = -\nabla f_i.
\]
Entretanto, reinicializações excessivas fazem com que o método se aproxime de um gradiente descendente, perdendo as vantagens de conjugação; reinicializações escassas, por outro lado, mantêm direções degradadas que acumulam erro numérico.

\subsection*{Resumo}

O método de Fletcher--Reeves funciona bem quando:
\begin{itemize}
    \item a função é quadrática ou quase quadrática;
    \item a matriz Hessiana é bem condicionada;
    \item há um \textit{line search} de boa qualidade.
\end{itemize}

Fora dessas condições, o método tende a apresentar perda de conjugação, sensibilidade numérica e degradação de desempenho.


\end{document}
